% AI-written, AI cross-reviewed; no human sign-off; not a claim that RH is proved
% Atlas unit D550. Local manuscript v3, item G.581, heading on printed PDF p. 302.
% Formalism.LFunctionChainParameters.LErr.large_x_eventually
% paper2/decomposition_tables.tex:4814–4819. Verbatim excerpt; SHA-256 7a72b3bd81a2279395b40207d5cd64f500e22b5b45dc627af3d2cfb4b453456e.
% Original TeX uses the macros and cross-reference keys of the source manuscript.

\begin{formalstatement}[D550]\label{decomp:D550}
For positive epsilon and fixed K,cHH,eK there exists x0 such that for every x at least x0 both log x and the constant KBig(K,cHH,eK) are at most x\textasciicircum{}epsilon.
\end{formalstatement}
{\small\raggedright\noindent\textsf{\bfseries Routine assembly}\quad \nolinkurl{Formalism.LFunctionChainParameters.LErr.large_x_eventually}\par\noindent \nolinkurl{Formalism/SizeWindowBrackets/ChainSlopeWindows/SubThreeMulLogTwoLe.lean:319--329}\par\noindent SHA-256 \nolinkurl{90f46cef73a8bbf9409ddb93252d5d50dd1093899eda3623171ad8cb61b9a738}\par\noindent Dominated 13\quad Disjoint 13\quad terminal.\par}
Intersects the eventual power bound for log with the eventual domination of a fixed constant by x\textasciicircum{}epsilon.
{\small\raggedright Local proof uses \nolinkurl{isLittleO_log_rpow_atTop}, \nolinkurl{tendsto_rpow_atTop}.\par}
